\chapter{Datos generales, manejo en tierra y servicio}
\label{cap:mm-datos-servicio}

\opening{Versiones, levantamiento, pesaje, estacionamiento y almacenamiento, lubricantes y servicio de cajas: el mantenimiento de todos los días.}

\section{Versiones del R22}
\fuente{MM, capítulo 16}

\begin{advertencia}
Las diferencias de números de parte entre versiones están en el Catálogo Ilustrado de Partes (IPC) del R22.
\end{advertencia}

\begin{center}
\small
\begin{tabularx}{\linewidth}{@{}>{\raggedright\arraybackslash}p{0.15\linewidth}>{\raggedright\arraybackslash}X@{}}
\toprule
Versión & Características \\
\midrule
Standard & Aprobado 16-mar-1979. S/N 0002--0199 (excepto 0175). Lycoming O-320-A2B o -A2C de 150~hp reducido a 124~hp; avgas 80/87 mínimo; 1300~lb; estabilizador A044-1 a 1,8°--2,3° nariz arriba. \\
HP & S/N 0175, 0200--0255, 0257--0300, 0302--0349, 0352--0356. O-320-B2C de 160~hp reducido a 124~hp; 100LL o 91/96 mínimo; 1300~lb; estabilizador 1,8°--2,3° nariz arriba. \\
Alpha & Aprobado 12-oct-1983. S/N 0256, 0301, 0350, 0351, 0357--0500. O-320-B2C a 124~hp; bastidores inferiores extendidos; cono de cola más alto; batería trasera; 1370~lb; estabilizador 2,8°--3,3° nariz abajo; tanque auxiliar opcional desde S/N 0457. \\
Beta & Aprobado 12-ago-1985. S/N 0501--2570. O-320-B2C a 124~hp con 131~hp por 5~min en despegue; panel de siete instrumentos; radiador de aceite mayor; 1370~lb; estabilizador 2,8°--3,3° nariz abajo. Desde S/N 2115, unidad A569-5 con bloqueo de arranque (impide arrancar por debajo de 62\,\% de RPM rotor si el actuador no está desacoplado del todo; permite arrancar en vuelo por encima de 69\,\%). \\
Mariner & Aprobado 12-sep-1985. S/N 0364 y 0501--2570 elegibles (sufijo ``M''). Como el Beta con flotadores utilitarios y protección anticorrosión adicional; cono distinto con estabilizador 1,8°--2,3° nariz arriba; batería en la nariz con flotadores; estabilizador de flotador en lugar del patín de cola. Sin flotadores, VFR día o noche; con flotadores, solo VFR diurno. \\
Beta II & Aprobado 31-ene-1996. S/N 2571 y posteriores. Lycoming O-360-J2A reducido a 124~hp continuos (131~hp 5~min); asistencia de calefacción del carburador; gobernador estándar; A569-5 con bloqueo de arranque; 1370~lb; estabilizador 2,8°--3,3° nariz abajo. \\
Mariner II & Aprobado 31-ene-1996. S/N 2571--3414 elegibles (sufijo ``M''). Beta II con flotadores; mismas restricciones que el Mariner. \\
Entrenador por instrumentos & Alpha, Beta o Beta II con panel de 10 instrumentos; solo operaciones VMC. \\
Policial & Alpha, Beta o Beta II con reflector, radios policiales y alternador de 70~A. \\
\bottomrule
\end{tabularx}
\end{center}

\textbf{Datum}: 100~in por delante del eje del rotor principal. Las estaciones (fuselaje, cono de cola, línea de agua y línea de simetría) se miden en pulgadas redondeadas a la centésima. Dimensiones externas: capítulo~\ref{cap:generalidades}. Los tornillos Torx Plus B526 (llave T20 o 20IP) con arandela de nailon B527-08 pueden usarse en capós y paneles (los B526-66 del estabilizador no son intercambiables).

\section{Levantamiento con gatos y con grúa}
\fuente{MM, capítulo 17}

\begin{precaucion}
Use gatos mecánicos o hidráulicos con bloqueo (un gato hidráulico sin bloqueo puede perder presión y alterar mediciones o desestabilizar la aeronave). Nunca levante con viento. Trabaje sobre superficie limpia, plana, dura y sin líquidos. No accione los gatos desde debajo del helicóptero y manténgase lejos de los patines. Al bajar, los patines se abren al cargar el peso.
\end{precaucion}

\begin{itemize}
\item \textbf{Gatos}: (a) un taco perfilado forrado de espuma \emph{directamente} bajo el cortafuego vertical, con un gato en cada extremo y un soporte bajo la cola; o (b) un gato bajo cada extremo del tubo transversal trasero, 1~in hacia adentro de los codos, y bajo la bola de manejo (que debe deslizar libremente). Accionar lenta y simultáneamente, una persona por gato; al menos 4~in sobre el suelo para retirar el tren.
\item \textbf{Grúa}: capacidad mínima \textbf{1000~lb}; cubo del rotor justo bajo el gancho. Con el accesorio MT234-1 (bujes de nailon en los agujeros de aligeramiento del cubo y pasadores elásticos) o con cuerda de nailon blanda de 1~in (carga de trabajo $\ge$1000~lb) en doble lazo por los agujeros del cubo. Estabilice guiando el patín de cola \emph{sin hacer fuerza} (es estructura secundaria).
\end{itemize}

\section{Nivelación, pesaje y lastre}
\fuente{MM, capítulo 18}

Nivele y pese siempre \textbf{sin viento}. Verifique el nivel de burbuja girándolo 180° sobre una superficie: la burbuja debe quedar igual.

\begin{itemize}
\item \textbf{Bastidor inferior derecho y tubo transversal trasero} (solo Standard y HP): nivel en la marca LEVEL HERE del bastidor y en el centro del tubo transversal.
\item \textbf{Cubo del rotor principal} (todos): nivel sobre el cubo, paralelo al perno de balanceo; con el perno alineado longitudinalmente se nivela longitudinalmente y con el perno transversal, lateralmente.
\item \textbf{Paneles de quilla} (todos): sobre el borde superior del panel derecho (longitudinal) y atravesado sobre ambos (lateral).
\end{itemize}

En cada caso se corrige con suplementos bajo los patines o gatos bajo el borde exterior del tubo transversal trasero (1~in hacia adentro de cada codo).

\subsection{Pesaje}

Repese cuando se modificaron el peso vacío y su CG y se duda de la exactitud de los cálculos. Lleve un registro continuo en el POH, sección 6.

\textbf{Preparación}: desabastecer; aceite del motor en servicio y cajas al centro del visor; aeronave limpia y seca; sin objetos sueltos; capós, paneles, puertas, mandos desmontables y POH instalados; lista de equipos (RF~134) coincidente con lo instalado.

\begin{pasos}
\item Eleve el helicóptero con grúa unos 30~cm, con las palas aproximadamente longitudinales y \textbf{levantadas de los topes de caída} para que el cubo bascule libre.
\item Colgando quieto, mida con un nivel de agua (sin burbujas) la diferencia de altura entre el eje de la caja de cola y el vientre de la cabina en el cortafuego vertical.
\item \textbf{CG longitudinal} = $114{,}47 - 0{,}315 \times$ (diferencia de altura, in).
\item Baje sobre balanzas de 1000~lb bajo cada patín, centradas unas 10~in por delante de los soportes de las ruedas de manejo, nivelado lateralmente. \textbf{Peso vacío} = derecha + izquierda $-$ tara.
\item \textbf{CG con combustible lleno y piloto solo mínimo}:
 \[ \mathrm{CG} = \frac{\mathrm{CG}_{\text{vacío}} \times W_{\text{vacío}} + M}{W_{\text{vacío}} + W} \]
 con $(M, W)$ = (28\,495; 303) vejigas con auxiliar; (22\,064; 240) vejiga sin auxiliar; (30\,180; 319) aluminio con auxiliar; (23\,011; 249) aluminio sin auxiliar.
\item Si queda detrás del límite trasero (POH, sección 2), lastre de nariz requerido $= \dfrac{(\mathrm{CG}_{\text{vacío}} - \text{límite}) \times W_{\text{vacío}} - 2358}{64{,}5}$~lb.
\item Corrija por el combustible no utilizable drenado (brazo 100,0~in): 10,2~lb con vejigas, 6,0~lb aluminio con auxiliar, 3,6~lb aluminio sin auxiliar; y por el lastre de nariz (brazo 37,5~in).
\item \textbf{CG lateral} $= \dfrac{\text{derecha} - \text{izquierda}}{\text{derecha} + \text{izquierda}} \times 37$~in.
\end{pasos}

\begin{precaucion}
El CG con combustible lleno y peso mínimo de piloto solo debe quedar en el límite trasero o por delante de él.
\end{precaucion}

\subsection{Lastre fijo}

\begin{center}
\begin{tabularx}{\linewidth}{@{}>{\raggedright\arraybackslash}Xccc@{}}
\toprule
Lastre & Máximo & Brazo long. & Brazo lat. \\
\midrule
De nariz (placas A941-2 de 0,25~lb y A941-3 de 1,00~lb) & 10,0~lb & 37,5~in & 0 \\
De empenaje (pesa A301-5) & 2,00~lb & 270,7~in & 2,6~in \\
\bottomrule
\end{tabularx}
\end{center}

El lastre de nariz va sobre la tapa A363-1 dentro de la consola inferior (agujeros de 0,198~in, broca \#8, usando la placa como plantilla). El de empenaje va bajo el soporte B902 o el estabilizador A044-1, sobre metal desnudo con recubrimiento de conversión e imprimación epoxi cromatada húmeda, con pernos NAS6604-44 (sin lastre, NAS6604-28). Todo cambio de lastre exige recalcular o repesar y actualizar el registro de peso y balance.

\section{Manejo en tierra}
\fuente{MM, capítulo 19}

\textbf{Ruedas de manejo}: van por fuera de los patines y son intercambiables. Inflado máximo \textbf{60~psi} (4~bar). Se inserta el eje completamente en el soporte y se tira de la manija en arco hacia la derecha hasta que quede trabada paralela al patín (con las ruedas trabadas, la manija derecha apunta adelante y la izquierda atrás). Sujete la manija con ambas manos: tiende a saltar.

\begin{advertencia}
Retire las ruedas de manejo antes del vuelo.
\end{advertencia}

Para mover el helicóptero: palas aproximadamente longitudinales y freno rotor aplicado; zona libre de objetos; una persona baja la cola hasta que los patines queden paralelos al suelo y guía tomando la caja de cola, la parte trasera del cono o la parte interior del estabilizador horizontal. \textbf{No tome} el protector del rotor de cola, los estabilizadores verticales, la parte exterior del horizontal, el rotor o sus controles, ni el patín de cola. Una segunda persona puede empujar un tubo vertical trasero o la nariz. Existen carros de manejo MT900-1 y MT950-1 para una sola persona. Los Mariner con flotadores usan un procedimiento con tacos y ruedas bajo los patines.

\textbf{Remolque en tráiler}: no recomendado (suspensiones de tráileres grandes que trabajan mal con poca carga y resonancia del cono de cola). Si es inevitable: lastrar el tráiler, apoyar el cono acolchado, desmontar las palas o apoyarlas a unos 5~ft de las puntas para descargar los topes, inmovilizar el rotor de cola, proteger parabrisas y equipos frágiles y, al terminar, hacer una \textbf{inspección de 100 horas}.

\section{Estacionamiento y almacenamiento}
\fuente{MM, capítulo 20}

\textbf{Estacionamiento}: cíclico neutro y colectivo abajo con fricción; palas aproximadamente longitudinales con freno rotor; amarres de palas con viento. \textbf{Nunca baje una pala tirando de ella: empuje hacia arriba la opuesta.}

\textbf{Viento fuerte o chorro de otras aeronaves}: hangar o zona segura ante tormentas; nariz hacia el viento; palas ligeramente desviadas de la posición longitudinal para que la trasera no golpee el cono; en hangar, válvula de combustible cerrada. En condiciones extremas, a criterio del operador, se desmontan palas y rotor de cola y se retiran o atan las bieletas (marcando ``X'' y ``O'').

\begin{precaucion}
Amarres de palas MT290-1 (letras blancas hacia abajo; la pala delantera al poste central del parabrisas y la trasera alrededor del cono delante de la luz anticolisión): \textbf{tensión máxima 5~lb}. Apretarlos de más daña las palas.
\end{precaucion}

\textbf{Almacenamiento de más de 30 días}:
\begin{pasos}
\item Registre la fecha.
\item Retire la batería (negativo primero) y mantenga según su fabricante.
\item Preserve el motor según Lycoming SL~L180.
\item Llene ambas cajas al centro del visor con el aceite correcto.
\item Cargue combustible: lubrica la vejiga e impide que se seque o agriete (recargue tras las pruebas de motor).
\item Limpie y seque; repinte o encere las palas.
\item Preferentemente en hangar; aplique un compuesto anticorrosivo no secante al eje del embrague A166 junto a los sellos de la polea superior.
\item Ventile la cabina y cierre las puertas de acceso. Inspeccione periódicamente por corrosión (más a menudo en ambientes húmedos).
\end{pasos}

\textbf{Calcomanías} (capítulo 21 del MM; ubicaciones en el IPC): se retiran despegando desde las esquinas a mano o con cuña plástica sin dañar pintura; se aplican sobre superficie limpia y seca (acetona o agua jabonosa; cuidado con la acetona cerca de la serigrafía de consola y panel de disyuntores).

\section{Servicio de las cajas de transmisión}
\fuente{MM, capítulo 22}

Evalúe el nivel con la aeronave nivelada y busque fugas cuando haga falta completar.

\begin{itemize}
\item \textbf{Drenaje de la caja principal}: calentar unos 5~min a 70--75\,\% de RPM; desconectar y retirar el detector de virutas; insertar el drenaje MT054-1 (MT053-1 para detector roscado) con la manguera por la borda y acuñarlo contra el cortafuego para abrir la válvula.
\item \textbf{Llenado}: hasta el centro del visor; girar el rotor a mano varias vueltas y bajar varias veces la caja de cola; verificar y corregir; tapón con torque especial; 5~min a 70--75\,\% y nueva verificación.
\item \textbf{Caja de cola}: se drena retirando el detector de virutas; se llena igual por el tapón; verificar la luz TR CHIP con el botón de prueba.
\end{itemize}

\subsection{Aceite de las cajas}

\begin{center}
\small
\begin{tabularx}{\linewidth}{@{}>{\raggedright\arraybackslash}p{0.2\linewidth}>{\raggedright\arraybackslash}X@{}}
\toprule
Aceite & Aplicación \\
\midrule
A257-2 (azul o azul verdoso) & Entregado en S/N anteriores a 4825. Todas las cajas (incluso nuevas) pueden usarlo. Calcomanía F654-14. \\
A257-22 (ámbar semitransparente, sintético) & Entregado desde S/N 4825. Caja principal A006 con piñón A146-1 rev.~Q o posterior; caja de cola con visor B563-4 (con filtro UV; no todas pueden adaptarse). Calcomanía F654-50. \\
\bottomrule
\end{tabularx}
\end{center}

\begin{precaucion}
No mezcle A257-22 con A257-2. Use el mismo aceite en ambas cajas, de acuerdo con la calcomanía de la pared trasera del tanque auxiliar (cámbiela si cambia de aceite). Para pasar de A257-22 a A257-2, o ante una mezcla accidental, drene y lave la caja.
\end{precaucion}

\textbf{Drenaje y lavado} (cajas principal y de cola): drenar en caliente; llenar con aceite mineral SAE~30, 40, 50 o 20W50; 5~min de marcha a 70--75\,\%; drenar; limpiar detector, alojamiento y visor; llenar con el aceite correcto y probar la luz de virutas.

\textbf{Detectores de virutas}: limpiar con cepillo de dientes y solvente aprobado; quitar residuos con aire comprimido o cinta adhesiva, \textbf{nunca con un imán}; secar e inspeccionar. Antes de instalar, conectar y tocar el bastidor con el imán: debe encenderse la luz correspondiente; después de instalar, verificar con el botón de prueba.

\section{Lubricación del embrague y de los rodamientos del actuador}

\subsection{Lubricante del embrague (sprag)}

En embragues con retén delantero A168-4: con los pernos de los yugos verticales y el freno rotor aplicado, se instalan racores MT147-2 en lugar de los tornillos B289-3, se drena y se lava la cámara del sprag con \textbf{A257-4} hasta que salga claramente rojo. Todo el lubricante drenado se cuela por un \textbf{filtro de pintura de 180--200~micras}: que el fluido salga oscuro y con brillo fino metálico es normal, pero \textbf{si el filtro retiene partículas metálicas, el embrague se devuelve a RHC} o a un centro autorizado. Luego se llena desde abajo hasta que no haya burbujas y se verifica que el nivel toque las roscas. Evite contaminar correas y poleas (limpiar con agua jabonosa y enjuagar). En embragues sin tornillos B289-3, se reemplaza el sello trasero (\S\,7.213; capítulo~\ref{cap:mm-transmision}).

\subsection{Rodamientos superior e inferior del actuador}

El rodamiento superior está en el eje del embrague y el inferior en el eje del ventilador. \textbf{Su falla en vuelo puede causar la pérdida de potencia al rotor y un accidente grave.} Justo antes de fallar, la luz del embrague puede parpadear constantemente (encendidos de menos de 1~s) o encenderse más tiempo de lo habitual; no confundir con el retensado normal (hasta 8~s). No vuelva a volar hasta determinar la causa. Inspecciónelos cuando se sospeche un problema o al retirar el ventilador.

\begin{itemize}
\item \textbf{Superior (A184)}: con el kit KI-115 (jeringa y sonda), inyecte \textbf{2--3~g de grasa A257-12} (2,5~g llenan 0,5~in de una jeringa de 0,63~in de diámetro interior) bajo el labio del sello trasero limpio, sin dañarlo; reponga el labio con la sonda; 2~min a 104\,\% y limpie la grasa expulsada. \textbf{No sustituya el rodamiento}: tiene holgura interna modificada para el R22.
\item \textbf{Inferior}: los A181-1 y A181-3 son obsoletos y deben retirarse de servicio. Con el ventilador fuera y el actuador desconectado, gire el alojamiento: sin aspereza ni raspado, \textbf{juego axial máx.\ 0,010~in}, sin daño en sellos ni por calor. Para lubricar se retira el bloque con la pista exterior partida, se limpian e inspeccionan las piezas y se llenan con A257-12 las aberturas a ambos lados de la jaula (no los espacios entre rodillos). Tornillos MS24694-S61 con torque estándar. \textbf{No se admite movimiento ni \emph{fretting} entre la pista interior y el eje del ventilador}; ante indicios de falla se reemplaza el conjunto eje y rodamiento A007-3/-5.
\end{itemize}

\section{Combustible: carga y descarga}

\begin{advertencia}
Cargue y descargue en lugar bien ventilado. Prohibido fumar a menos de 100~ft.
\end{advertencia}

\textbf{Carga}: batería apagada; conectar a tierra helicóptero y vehículo, y el vehículo al helicóptero; verificar el grado de combustible (POH, sección 2, incluido el uso de anticongelante). Con tanques de vejiga, la tapa A689-6 debe ser rev.~J o posterior. \textbf{Descarga}: batería y válvula cerradas; recipiente conectado a tierra bajo el carburador; desconectar la manguera del carburador y abrir la válvula (puede aprovecharse para probar el interruptor de bajo combustible); reconectar con torque especial y marca.

\section{Cabina, puertas y parabrisas}
\fuente{MM, capítulos 15 y 25}

\begin{itemize}
\item Los asientos no son regulables; se entrega un cojín para adelantar a pilotos de baja estatura. Cada asiento tiene cinturón con arnés de hombro de carrete inercial y un compartimiento de equipaje debajo.
\item \textbf{Hebillas F628-8} giradas \textbf{20° hacia adelante} respecto de la vertical; pernos de anclajes y guías con torque estándar. Placa A798-1 de la válvula de combustible a 16~in-lb. Respaldo de cojín centrado 1,00 $\pm$ 0,10~in sobre el borde del panel.
\item Alfombras y aislamiento se pegan con adhesivo B270-8 (no se recomienda reutilizar el aislamiento).
\item \textbf{ELT Kannad}: registre la baliza al comprarla y ante cambios de contacto, propietario o matrícula (cospas-sarsat.org); la llave (\emph{dongle}) debe programarse con los datos de la aeronave.
\item \textbf{Extintor}: temperatura de operación y almacenamiento $-$40~°F a 120~°F ($-$40~°C a 49~°C); en climas calurosos al sol, retírelo. Inspección: sin daños ni corrosión, boquilla libre, precinto intacto, calcomanía de registro legible; pese con balanza de 1~g y retírelo si está por debajo del mínimo. El B381-1 se retira al final del \textbf{12.º año} de fabricación. Tornillos del soporte en la nariz con 0,00--0,06~in de rosca expuesta.
\item \textbf{Puertas}: pueden retirarlas mecánicos o pilotos. \textbf{Cada uno de los dos pasadores de bisagra debe llevar chaveta}; sin ella la puerta puede desprenderse en vuelo. Ajuste el peso y balance al quitar o poner puertas.
\end{itemize}

\subsection{Parabrisas}

Dos mitades en marcos de aluminio con un arco tubular central. Existen parabrisas \textbf{acrílicos} (estándar), entre retenedores con sellador, y \textbf{de policarbonato} resistentes a impactos, fijados además con flejes de acero inoxidable G367. El policarbonato es más resistente y flexible pero más sensible a la radiación UV y se raya con facilidad: limpie con cuidado, evite solventes incompatibles (fragilización o cuarteado) y use funda o hangar.

\textbf{Límites de inspección} (junto a los retenedores):
\begin{itemize}
\item Cuarteado (fisuras que no atraviesan): total no mayor de 2~in de largo, ninguna de más de 1~in, al menos 1~in libre entre zonas y no más de 1/2~in desde el borde del retenedor. Reemplace si se nota con la uña o perjudica la visión.
\item Fisuras no paralelas al borde y de más de 0,25~in: reemplazar. Las menores de 0,25~in pueden detenerse con un agujero de 3/32~in.
\item En policarbonato, revise también junto a fijaciones y flejes, y decoloración o turbidez significativas.
\end{itemize}

Para cambiarlo se protege con papel o película, se instalan las puertas para soportar la estructura, se recorta el nuevo con sierra de cinta (acrílico: $\ge$24 dientes por pulgada; policarbonato: ángulo de despeje 20°--30°, ataque 0°--5°, 600--1000~m/min, paso 1,5--4~mm) sobre cartón, se alisan los bordes con lija 100 y se arma con sellador B270-26 y cinta de enmascarar para el exceso (limpiar con lacolene). En el policarbonato, los flejes deben quedar perpendiculares a los retenedores dentro de 2° y los agujeros se agrandan a 0,169--0,175~in.
